\section{Normative scientific requirements}
\label{sec:requirements}

The words ``must'', ``must not'', ``required'', and ``invalid'' are normative.
These requirements define scientific conformance without prescribing a file,
class, function, language, or storage layout.

\subsection{Scope, state, and input typing}

\begin{description}[style=nextline]
\item[\Req{001}: measured-channel scope.]
  The v0.1 operation must accept only the ordinary \code{xs} measured stream.
  A physically distinct measured stream must be rejected as unsupported and
  must not inherit the \code{xs} factor, unit, response, or uncertainty label.

\item[\Req{002}: target-unit scope.]
  The only calibrated target is top-of-atmosphere point-source-peak
  \(\mjybeam\) at \(X_{\rm ref}=0\). Its photometric-convention state must be
  explicit. While Q05 is open, complete monochromatic mJy meaning and
  cross-array spectral comparison must remain unavailable. An unsupported
  unit or photometric meaning must yield no calibrated SCI-CAL output, not a
  relabeled numeric array.

\item[\Req{003}: complete signal declaration.]
  Each admitted signal must declare its acquisition identity, pre-CAL unit,
  sample and detector shape, time/support, upstream eligibility, and any
  authorized offset/baseline convention. It must resolve the approved physical
  observable, sign, preprocessing, normalization, stream scope, and valid
  multiplicative loading/Tune/signal domain. Missing, incompatible, or
  authority-unresolved declarations invalidate the corresponding physical
  interpretation or calibrated-unit output exactly as Q01 records.

\item[\Req{004}: explicit semantic states.]
  Missing, disabled, automatic, unavailable, invalid, engineering-only, and
  science-qualification-eligible states must be represented explicitly. None
  may be encoded solely as a zero, unity, negative, non-finite, stale, or
  undocumented sentinel.

\item[\Req{005}: one-way resolution.]
  Requested, effective, observation-resolved, and realized calibration states
  must remain distinct and flow one way. Realized or compatibility state must
  not mutate the accepted request, and a new observation must replace its
  resolved state in full. The owner-decision snapshot is immutable input to
  resolution; candidate or realized state must not mark an open decision
  closed.
\end{description}

\subsection{Occurrence-scoped identity and association}

\begin{description}[style=nextline]
\item[\Req{006}: acquisition key.]
  Every selected detector occurrence must have an observation-local key
  containing observation/Tune context, network or interface, and the
  network-local channel/tone slot. A global column is only a locator.

\item[\Req{007}: binding evidence.]
  A sample column may be bound to an acquisition occurrence and selected APT
  row only by an explicit keyed relation or a verified ordered-row contract
  that records scope, order definition, cardinality, and proof. Row position
  without that evidence is invalid.

\item[\Req{008}: unique coverage.]
  On the declared selected support, binding and target-to-source association
  must be single-valued and complete. Missing, duplicate, ambiguous, or
  contradictory keys must invalidate exactly the affected detector or the
  predeclared required support; no arbitrary choice is permitted.

\item[\Req{009}: immutable source and selected-child APT identity.]
  The source APT and selected observation-specific child APT must each be
  identified as one immutable artifact by resolvable artifact identity and
  digest. Each used row must carry its artifact-local key and occurrence
  context. Parent and child remain distinct even for an identity
  transformation or numerically unchanged value.

\item[\Req{010}: association evidence.]
  Each admitted target-to-source edge must name both occurrence-scoped
  endpoints, matcher method/version, disposition, and available quality
  evidence. Only an admitted unique match may calibrate; abstained, ambiguous,
  rejected, or missing edges must not. The record must identify TolProj as the
  association/child-transform authority and must not assign source-factor
  meaning to the matcher or consumer.

\item[\Req{011}: layered identity.]
  Acquisition, measured-APT, association, and design identities, as well as
  occurrence, semantic-content, and byte-transport identities, must remain
  distinct. A design match is not required for ordinary measured Beammap
  calibration, and a design-ID change alone must not change a valid result.

\item[\Req{012}: array/network consistency and lifetime.]
  Network-to-array membership and target/source network context must be
  checked as declared constraints. Reordering is harmless only when the
  binding proof remains valid. No identity or row state may leak from another
  observation.
\end{description}

\subsection{Factor authority, causality, and reconstructibility}

\begin{description}[style=nextline]
\item[\Req{013}: absolute factor.]
  The sole absolute detector factor must be the selected child-APT row's finite,
  nonzero \code{flxscale}, with direction \(\mjybeam/U_x\), detector recipient,
  validity, reference plane, uncertainty status, and selected-artifact lineage
  declared. Its source-calibration record must resolve producer identity,
  source APT/row, calibrator model and epoch, source-atmosphere treatment,
  amplitude estimator, beam/template normalization, pointing treatment,
  spectral convention, factor direction, and uncertainty/covariance status.
  Reciprocal, sign, or missing generating meaning must never be inferred.

\item[\Req{014}: pointing-correction disposition.]
  The selected lineage must name the TolProj association and approved child
  transformation, its version, parent and child occurrences,
  transfer-validity domain, and retained systematics. It must classify any
  pointing-derived correction
  as \code{not_applicable}, \code{embodied_once}, or
  \code{required_but_unresolved}. An embodied correction must name its parent
  and event and must not be applied again; an unresolved required correction
  prohibits calibrated output.

\item[\Req{015}: once-only application.]
  Absolute \code{flxscale}, target atmosphere correction, and the explicit
  target-unit identity must each have realized application count one on every
  valid occurrence. Parent factors and already embodied events must have
  runtime count zero.

\item[\Req{016}: canonical multiplier.]
  The signal must obey Equations~\eqref{eq:multiplier} and
  \eqref{eq:signal} on exactly \(\valid\). Scalar multiplication order may
  change without changing values, but the causal lineage and factor-instance
  identities must be preserved.

\item[\Req{017}: responsivity boundary.]
  \code{responsivity} must not be presented or applied as an independent
  absolute \(\mjybeam\) factor. Any use must have a distinct relative role,
  unit, normalization/recipient, support, and lineage, and must not duplicate
  information already embodied in \code{flxscale}.

\item[\Req{018}: sensitivity boundary.]
  Beammap \code{sens} must not multiply the calibrated signal or masquerade as
  atmosphere correction or total uncertainty. If used for an approximate
  weight, that role, units, estimator identity, approximation, support, and
  exclusions must be explicit.

\item[\Req{019}: compatibility factor boundary.]
  An \code{fcf} or other compatibility aggregate may be admitted only after
  exact decomposition into the canonical role, unit, direction, constituent
  instance identities, and exclusions. An opaque aggregate has no calibration
  authority.

\item[\Req{020}: factor and package lineage.]
  Every realized factor must record definition, role, unit, direction,
  recipient/support, reference plane, application stage, validity,
  uncertainty status, and parent lineage. One canonical record per coherent
  reduction package must distinguish the Beammap/source-APT producer, TolProj
  association/child transformer, TolTECA delivery, and SCI-CAL consumer; bind
  the Q01--Q09 snapshot, photometric-convention state, and pipeline boundary;
  and reconstruct Equation~\eqref{eq:reconstruction}. Product-level links may
  resolve it without duplicating the complete APT.
\end{description}

\subsection{Atmosphere operator, support, and quality}

\begin{description}[style=nextline]
\item[\Req{021}: opacity and airmass meaning.]
  \(\tau_{225}\) must be finite nonnegative zenith optical depth with source
  and time support. \(X\) must be the full eligible-sample airmass with model
  identity/domain. Equation~\eqref{eq:los-coordinate} and
  \(X_{\rm ref}=0\) must be used; a zenith value alone is not a line-of-sight
  correction.

\item[\Req{022}: temporal support.]
  Every calibrated sample must be covered by an exact opacity state or a
  source-authorized interpolation between valid bracketing states. Method,
  version, endpoints, validity interval, and any maximum gap must be recorded.
  Unbracketed interpolation, endpoint extrapolation, and inherited prior-
  observation opacity are prohibited.

\item[\Req{023}: complete operator identity.]
  A numeric atmosphere evaluation requires the exact operator name, content
  digest, array assignment, ordered coordinate/ordinate nodes with units,
  ordinate orientation, interpolation rule, numeric support, generating model
  provenance, and passband identity when used. A name without these facts is
  insufficient. Operator completeness closes only Q06; it does not establish
  source-factor derivation, transfer validity, broadband compatibility, or
  end-to-end ordering.

\item[\Req{024}: structural evaluation.]
  The retained operator must obey Equations~\eqref{eq:operator-nodes} through
  \eqref{eq:operator-orientation}: analytic zero anchor, interpolation of the
  authoritative ordinate, exact node recovery, and continuity at seams.

\item[\Req{025}: physical-direction invariants.]
  On declared positive-opacity support, transmission and correction must be
  finite and positive, transmission must be strictly below unity and
  decreasing, and correction must be strictly above unity and increasing.
  A finite unity plateau, sign change, non-finite value, or direction reversal
  is invalid.

\item[\Req{026}: domain intersection.]
  A calibrated sample must satisfy Equation~\eqref{eq:atmos-domain}. No
  opacity, elevation/airmass, line-of-sight coordinate, or time request may be
  silently clamped, extrapolated, or replaced by a legacy selector/fallback.

\item[\Req{027}: science-qualification eligibility policy.]
  A coherent segment may receive the
  \code{science-qualification-eligible} structural-policy class only when all
  its admitted opacity support lies in \(0\leq\tau_{225}\leq0.15\) and every
  other structural admission condition holds. This class records eligibility
  to seek separately accepted qualification evidence; it is not an
  atmosphere-fidelity, observational-performance, \code{science-qualified},
  or \code{calibrated-science} claim.

\item[\Req{028}]
  \Term{Engineering-opacity restriction.} If any part of an unsplit coherent segment has
  \(0.15<\tau_{225}\leq0.25\), its class must be
  \code{engineering-only} \code{unavailable} and v0.1 must emit no calibrated
  SCI-CAL signal. It must not extrapolate, fall back, or preserve a calibrated
  unit label.

\item[\Req{029}: invalid and outside-supported opacity.]
  Negative or non-finite opacity is invalid input. Opacity above 0.25, absent
  opacity, invalid airmass/elevation meaning, or a request beyond any operator
  node/time/model support is outside supported calibration and must yield no
  calibrated output.

\item[\Req{030}: coherent quality state.]
  One predeclared observation or segment must receive one opacity quality
  class. It must take the most restrictive applicable state across its
  support; sample-by-sample class switching is prohibited. A separately
  declared split creates separate lineage-bearing segments.

\item[\Req{031}: passband limits.]
  Where the operator depends on a passband, it must bind the exact selected
  TolTECA v1 passband-set identity and record all known unknowns in
  Section~\ref{sec:definitions}. No detector/network/telescope measurement,
  normalization, uncertainty, covariance, or photon/energy semantics may be
  inferred from the array-named curves. A separate compatibility disposition
  must connect the source-factor, target-atmosphere, and reported-output
  photometric conventions; equal passband identity alone is not proof.
\end{description}

\subsection{Conditional and nuisance uncertainty}

\begin{description}[style=nextline]
\item[\Req{032}: full conditional covariance.]
  A supplied conditional covariance must transform by
  Equation~\eqref{eq:conditional-covariance} on the identical ordered valid
  support. Existing off-diagonal correlations must be preserved.

\item[\Req{033}: scalar variance and weight.]
  Supplied diagonal conditional variance or inverse variance must transform by
  Equation~\eqref{eq:variance-weight}, with output units
  \((\mjybeam)^2\) or \((\mjybeam)^{-2}\). Weight is not support, validity,
  hits, response, confidence, or total uncertainty.

\item[\Req{034}: unavailable conditional uncertainty.]
  If no admissible variance, weight, or covariance exists for a valid signal,
  conditional uncertainty must be marked unavailable on that same support.
  It must not be synthesized as zero variance, infinite weight, or a value
  derived from \code{sens} without a separately approved estimator.

\item[\Req{035}: nuisance completeness ledger.]
  The lineage must carry an entry for detector \code{flxscale}, common
  absolute calibrator scale, any TolProj pointing-derived correction,
  WVR/atmosphere model, Beammap \code{sens} when used for approximate weights,
  and beam/template response. Each entry must be \code{quantified},
  \code{not_applicable} with justification, or \code{unavailable}, and name
  detector, array, observation, cohort, or global correlation scope when
  known.

\item[\Req{036}: correlation preservation.]
  Nuisance covariance must preserve declared common and cross terms. Shared
  calibrator, array, opacity, pointing, passband, or beam terms must not be
  treated as independent per sample or counted down solely with sample or
  detector number.

\item[\Req{037}: total-uncertainty claim.]
  Total calibrated uncertainty or significance is available only when every
  required nuisance category and material cross-covariance is quantified or
  explicitly shown not applicable. Otherwise it must be labeled unavailable,
  even when conditional weight exists.

\item[\Req{038}: propagation regime.]
  Linear nuisance propagation may be claimed only inside one differentiable
  operator interval with sufficiently small approximately symmetric
  uncertainties and justified dependence assumptions. Seam-spanning, large,
  asymmetric, selection, association, or same-data fitted-state uncertainty
  requires nonlinear or joint propagation; missing cross terms may not be set
  to zero without justification.

\item[\Req{039}: shared realized operator.]
  Signal, conditional covariance/weights, admitted synthetic injections,
  conditional noise realizations, and genuine Jacobian companions must use
  the identical realized multiplier, valid support, and factor lineage. A
  different measured stream must not be substituted as a companion. This
  local equality does not establish a unique end-to-end response while the
  Q02 pipeline boundary/order decision is open.
\end{description}

\subsection{Response basis, validity, and claims}

\begin{description}[style=nextline]
\item[\Req{040}: originating response basis.]
  Every calibrated result must resolve its producer-owned Beammap source,
  source APT, selected child APT, and per-
  detector beam/template occurrence. Available elliptical major/minor axes,
  position angle, units, coordinate frame, epoch/support, fit identity, and
  uncertainty status must be retained.

\item[\Req{041}: realized response distinction.]
  Originating Beammap response and realized mapmaker/kernel/filter response
  must remain separate provenance objects. Circularization must be explicitly
  labeled with its rule and cannot silently replace available elliptical
  information.

\item[\Req{042}: point-source-peak condition.]
  A downstream result may claim literal point-source-peak \(\mjybeam\) only
  when its realized response satisfies Equation~\eqref{eq:peak-response} with
  \(\rho_L=1\) for the declared beam/template and support, or when an explicit
  response renormalization establishes and records that condition.

\item[\Req{043}: response-changing operations.]
  A kernel, filter, or map operation that changes peak response must propagate
  the response basis and factor. Without a separately established
  renormalization it must withhold the literal peak-calibration claim; SCI-CAL
  itself must not certify the downstream response.

\item[\Req{044}: excluded photometric meanings.]
  V0.1 must not emit or authorize \(\mathrm{MJy/sr}\), \(\mathrm{Jy/pixel}\),
  Rayleigh-Jeans or thermodynamic temperature, extended-source calibration,
  integrated photometry, or a resolved-source recovery claim.

\item[\Req{045}: validity tuple and causes.]
  Every coherent result must publish the tuple in
  Equation~\eqref{eq:validity-tuple} and machine-distinguishable reason codes
  for missing identity, unsupported stream/unit, invalid or out-of-domain
  atmosphere, intentionally disabled calibration, unavailable conditional or
  nuisance uncertainty, incomplete response basis, and any required output
  failure. The associated owner-decision snapshot must make every Q01--Q09
  limitation machine-distinguishable without converting an open decision into
  a numeric sentinel.

\item[\Req{046}: truthful no-output states.]
  Disabled/uncalibrated, invalid, engineering-only unavailable, outside-
  supported calibration, and uncertainty-unavailable states must remain
  distinguishable. A failure in a required factor or identity must not emit a
  calibrated number; a missing non-signal uncertainty may leave the signal
  valid only with the corresponding claim explicitly unavailable. Missing
  owner authority follows the exact decision-ledger consequence: it blocks
  only the affected output or interpretation and is never silently promoted.

\item[\Req{047}: canonical lineage record.]
  One resolvable canonical record per coherent reduction package must include
  raw observation identity, source and selected-child APTs and digests,
  acquisition and row occurrences, association edge, approved child
  transformation, delivery identity, factor definition and generating record,
  atmosphere/operator/passband identities and support, reference plane,
  target unit, photometric-convention state, pipeline boundary/order state,
  Q01--Q09 snapshot, state tuple, uncertainty ledger, and response basis.

\item[\Req{048}: compact product links.]
  Every calibrated product must contain links and occurrence keys sufficient
  to resolve its canonical lineage, exact source APT, selected child APT, and
  governing owner-decision snapshot. It need not and should not duplicate the
  entire APT or dense per-detector lineage table in every product.

\item[\Req{049}: separated claim status.]
  Structural correctness, atmosphere representation fidelity, relative
  repeatability, and absolute flux performance must be reported as separate
  claims with separate evidence and support. A pass in one layer must not be
  promoted to another. In particular,
  \code{science-qualification-eligible} must not be promoted to an achieved
  \code{science-qualified} or \code{calibrated-science} claim. Unless the owner
  later defines a different explicit threshold, either achieved claim requires
  separately accepted atmosphere-representation-fidelity and
  observational-performance evidence over the same declared identity and
  support in addition to structural correctness. Input/factor physical
  meaning, calibrator-to-target transfer, broadband photometric meaning, and
  end-to-end response are also separate claim-specific dispositions governed
  by Q01--Q05.

\item[\Req{050}: falsifiability before qualification.]
  Any conformance or performance claim must identify evidence capable of
  falsifying the relevant equations, assumptions, edge cases, identity
  invariants, uncertainty correlations, and response condition. Skipped
  required evidence is unavailable, not a pass; the provisional one-percent,
  five-percent, and five-to-ten-percent targets are not guarantees. The
  evidence package must show that resolving Q06 alone does not promote any
  Q01--Q05 or Q07--Q09 claim.
\end{description}
